\documentclass[12pt,a4paper]{ctexart}
\usepackage[left=2.5cm,right=2.5cm,top=2.3cm,bottom=2.3cm]{geometry}
\usepackage{xcolor}
\usepackage{tcolorbox}
\usepackage{titlesec}
\usepackage{fancyhdr}
\usepackage{enumitem}
\usepackage{tabularx}
\usepackage{amsmath,amssymb}
\usepackage{hyperref}
\tcbuselibrary{skins,breakable}

\definecolor{MainColor}{HTML}{2A6F73}
\definecolor{SecondColor}{HTML}{3CAEA3}
\definecolor{LightMain}{HTML}{EAF7F6}
\definecolor{TitleDark}{HTML}{1F3A40}
\definecolor{TextGray}{HTML}{555555}
\definecolor{BorderGray}{HTML}{CBD5D6}
\definecolor{WarnColor}{HTML}{F2B35D}
\definecolor{ErrorColor}{HTML}{C94C4C}
\definecolor{DefineColor}{HTML}{6C63B7}

\linespread{1.25}
\setlength{\parindent}{2em}
\setlength{\parskip}{0.25em}
\hypersetup{colorlinks=true,linkcolor=MainColor,urlcolor=MainColor}
\pagestyle{fancy}
\fancyhf{}
\fancyfoot[L]{\textcolor{MainColor}{机器学习笔记}}
\fancyfoot[R]{\textcolor{MainColor}{\thepage}}
\renewcommand{\headrulewidth}{0pt}
\renewcommand{\footrulewidth}{0pt}

\newcommand{\topbar}[2]{\noindent{\small\textcolor{MainColor}{#1}}\hfill{\small\bfseries\textcolor{TitleDark}{#2}}\par\vspace{0.35em}{\color{BorderGray}\hrule}\vspace{1.2em}}
\newcommand{\notetitlecard}[6]{
\begin{tcolorbox}[enhanced,colback=white,colframe=BorderGray,boxrule=0.6pt,arc=2mm,left=12pt,right=12pt,top=10pt,bottom=10pt,borderline west={4pt}{0pt}{SecondColor}]
\begin{tabularx}{\linewidth}{X r}
{\small\textcolor{TextGray}{#1}} & {\small\textcolor{TextGray}{#2}} \\[0.9em]
\multicolumn{2}{l}{{\Huge\bfseries\textcolor{MainColor}{#3}}} \\[0.45em]
\multicolumn{2}{l}{{\large\textcolor{SecondColor!90!black}{#4}}} \\[0.45em]
\multicolumn{2}{l}{{\small\textcolor{TextGray}{课程：#5 \qquad 作者：#6}}}
\end{tabularx}\end{tcolorbox}\vspace{1em}}
\newtcolorbox{summarybox}[1][本节摘要]{enhanced,breakable,colback=LightMain,colframe=SecondColor,colbacktitle=SecondColor,coltitle=white,title={#1},fonttitle=\bfseries,boxrule=0.6pt,arc=2mm,left=12pt,right=12pt,top=10pt,bottom=10pt}
\newtcolorbox{definitionbox}[1][定义]{enhanced,breakable,colback=DefineColor!6,colframe=DefineColor,colbacktitle=DefineColor,coltitle=white,title={#1},fonttitle=\bfseries,boxrule=0.6pt,arc=2mm,left=12pt,right=12pt,top=10pt,bottom=10pt}
\newtcolorbox{keybox}[1][重点]{enhanced,breakable,colback=MainColor!5,colframe=MainColor,colbacktitle=MainColor,coltitle=white,title={#1},fonttitle=\bfseries,boxrule=0.6pt,arc=2mm,left=12pt,right=12pt,top=10pt,bottom=10pt}
\newtcolorbox{examplebox}[1][例题]{enhanced,breakable,colback=white,colframe=SecondColor,colbacktitle=SecondColor!85!black,coltitle=white,title={#1},fonttitle=\bfseries,boxrule=0.6pt,arc=2mm,left=12pt,right=12pt,top=10pt,bottom=10pt}
\newtcolorbox{mistakebox}[1][易错点]{enhanced,breakable,colback=ErrorColor!5,colframe=ErrorColor,colbacktitle=ErrorColor,coltitle=white,title={#1},fonttitle=\bfseries,boxrule=0.6pt,arc=2mm,left=12pt,right=12pt,top=10pt,bottom=10pt}
\newtcolorbox{tipbox}[1][提示]{enhanced,breakable,colback=WarnColor!10,colframe=WarnColor,colbacktitle=WarnColor,coltitle=white,title={#1},fonttitle=\bfseries,boxrule=0.6pt,arc=2mm,left=12pt,right=12pt,top=10pt,bottom=10pt}
\titleformat{\section}{\Large\bfseries\color{MainColor}}{\thesection}{0.8em}{}
\titleformat{\subsection}{\large\bfseries\color{TitleDark}}{\thesubsection}{0.8em}{}
\titleformat{\subsubsection}{\normalsize\bfseries\color{SecondColor!80!black}}{\thesubsubsection}{0.8em}{}
\setlist[itemize]{leftmargin=2.2em,itemsep=0.3em,topsep=0.4em}
\setlist[enumerate]{leftmargin=2.2em,itemsep=0.3em,topsep=0.4em}
\newcommand{\important}[1]{\textbf{\textcolor{MainColor}{#1}}}
\newcommand{\warning}[1]{\textbf{\textcolor{ErrorColor}{#1}}}
\newcommand{\concept}[1]{\textbf{\textcolor{DefineColor}{#1}}}

\begin{document}
\topbar{吴恩达机器学习}{学习笔记}
\notetitlecard{Study Notes Series}{2026 Spring}{吴恩达机器学习笔记(10)}{}{机器学习}{C++与草稿纸}

\begin{summarybox}
本周讨论大规模机器学习和图片文字识别。重点是随机梯度下降、小批量梯度下降、在线学习以及分布式数据处理，并用图片文字识别说明如何设计端到端的机器学习管道、滑动窗口和数据上限分析。
\end{summarybox}

\section{大规模数据集的学习}
\begin{keybox}[先判断是否真的需要更多数据]
大数据不是自动带来好效果。应先画学习曲线，判断模型是高偏差还是高方差；若验证误差随数据增加仍然很高，增加数据可能无效，需要修改特征或模型。只有在高方差场景下，扩大数据集通常才有明显收益。
\end{keybox}
\section{随机梯度下降}
\begin{definitionbox}[批量与随机梯度下降]
批量梯度下降每次使用全部 $m$ 个样本：
\[
\theta\leftarrow\theta-\alpha\frac1m\sum_{i=1}^{m}\nabla L_i(\theta).
\]
随机梯度下降（SGD）逐个样本更新：
\[
\theta\leftarrow\theta-\alpha\nabla L_i(\theta),
\qquad i=1,\ldots,m.
\]
SGD 每次更新便宜、可快速处理海量数据，但目标函数会有噪声，通常在最优点附近波动。
\end{definitionbox}
\begin{keybox}[小批量梯度下降]
将连续 $b$ 个样本组成一个小批量：
\[
\theta\leftarrow\theta-\alpha\frac1b\sum_{i=t}^{t+b-1}\nabla L_i(\theta).
\]
批量大小 $b$ 常取 $32$、$64$、$128$ 等，既能利用向量化和 GPU，又比全批量更新更快。每个 epoch 前可以随机打乱样本，减少固定顺序带来的偏差。
\end{keybox}
\begin{tipbox}[收敛与学习率]
观察每个 epoch 的平均代价而不是单个样本损失。若平均代价总体下降但有小幅波动，可能是 SGD 的正常噪声；若持续上升，检查学习率是否过大、特征是否缩放以及实现是否有错误。也可以逐渐衰减学习率，但衰减策略必须在验证集上评估。
\end{tipbox}

\section{在线学习}
\begin{definitionbox}[在线学习]
在线学习不断接收一个新样本 $(x,y)$，立刻用它更新参数，然后丢弃或存档该样本。它适合广告点击率、搜索排序或物流需求等数据分布随时间变化的场景。
\[
\theta\leftarrow\theta-\alpha\nabla L(\theta;x,y).
\]
\end{definitionbox}
\begin{mistakebox}[在线学习的风险]
在线数据可能包含垃圾样本、反馈回路或分布漂移。需要监测指标、保留离线验证集、限制异常更新，并定期用新旧数据重新评估模型。在线学习不是跳过数据清洗和模型验证。
\end{mistakebox}

\section{映射化简与数据并行}
\begin{definitionbox}[MapReduce 思路]
把大数据分到多个机器上：Map 阶段在各节点计算局部损失和梯度，Reduce 阶段把局部结果求和或平均，再执行一次参数更新。只要局部梯度的组合与全数据梯度等价，就可以在不把所有数据放到一台机器上的情况下训练模型。
\end{definitionbox}
\begin{tipbox}[分布式训练注意事项]
通信成本、节点故障、数据分片是否均衡和参数同步方式都会影响速度。异步更新吞吐较高但梯度可能过期，同步更新更稳定但要等待最慢节点。工程上需根据模型规模和网络条件选择方案。
\end{tipbox}

\section{应用实例：图片文字识别}
\subsection{问题描述和流程图}
\begin{keybox}[端到端管道]
图片文字识别通常拆为文字区域检测、字符分割或序列识别、字符分类和结果拼接等阶段。每个模块都有输入、输出和可测量指标，错误可能在任意阶段累积，因此不能只看最终准确率。
\end{keybox}
\subsection{滑动窗口}
\begin{definitionbox}[滑动窗口检测]
使用固定大小的窗口从左到右、从上到下扫描图像，对每个窗口预测“是否包含文字”。为了检测不同大小的文字，可以使用多个窗口尺寸或对图像做多尺度缩放。检测结果通常需要后处理合并重叠区域。
\end{definitionbox}
\subsection{获取数据与上限分析}
\begin{keybox}[数据和人工上限]
人工检查不同阶段的错误，估计每个模块距离理想性能还有多少空间。如果某一模块的上限已经接近最终指标，而另一模块仍有较大差距，应优先改进后者。合成数据、迁移学习和主动收集困难样本可以提高数据利用效率。
\end{keybox}
\begin{examplebox}[误差预算]
若文字检测准确率为 $95\%$、字符识别准确率为 $96\%$，最终正确率不会简单等于两者之和，而会受到错误传播影响。应分别测量“检测正确但识别错误”“检测遗漏”等类型，再决定增加标注、修改窗口还是更换分类器。
\end{examplebox}

\section{课程总结}
\begin{keybox}[从模型到系统]
本课程的主线是：线性回归和逻辑回归建立基本监督学习工具；神经网络和 SVM 扩展非线性表示；聚类、PCA 和异常检测处理无监督问题；推荐系统学习用户偏好；大规模学习则关注数据、计算和部署。实际项目中，算法只是系统的一部分，数据质量、验证设计、误差分析和工程约束同样决定最终效果。
\end{keybox}
\end{document}
